\documentclass[11pt]{article}
\usepackage[margin=1in]{geometry}
\usepackage{amsmath,amssymb,amsthm,mathtools}
\usepackage{booktabs,graphicx,microtype}
\usepackage[hidelinks]{hyperref}
\newtheorem{theorem}{Theorem}
\newtheorem{lemma}[theorem]{Lemma}
\newcommand{\R}{\mathbb R}
\newcommand{\rt}{\sqrt2}
\setlength{\emergencystretch}{2em}

\title{A period-eight Chebyshev factorization\\
for the zeros of a Jacobi-type pencil}
\author{Mathematical discovery run \texttt{fa\_banach\_001}}
\date{21 August 2026}

\begin{document}
\maketitle

\noindent\textbf{Status.}
Candidate full solution, independently rechecked and likely valid, pending
expert review.

\section{Source problem}

For the basic Jacobi-type pencil $\Theta_1$, Zagorodnyuk defines degree-$n$
polynomials $p_n$ and proves the representation
\begin{equation}\label{eq:source}
 q_n(t):=p_n(\rt t-1)
 =T_n(t)+tU_{n-1}(t)
 -\frac12\frac{U_{n-1}(t)-U_{n-1}(-1/\rt)}{t+1/\rt},
\end{equation}
initially away from $t=-1/\rt$ and hence everywhere by polynomial
continuation.  Here $T_j,U_j$ are the Chebyshev polynomials of the first and
second kind and $U_{-1}=0$.  The source asks for the zero distribution and
conjectures that all zeros are real \cite[formula (8) and PDF p.~6]{Z}.

\begin{center}
\includegraphics[width=0.96\textwidth]{figures/open_problem_crop.png}
\end{center}

\section{Idea of the proof}

The divided difference in \eqref{eq:source} hides a periodic factorization.
Multiplication by $2(t+1/\rt)$ turns $q_n$ into
$U_{n+1}+\rt U_n+c_n$, where
$c_n=\rt\sin(3\pi n/4)$ depends only on $n$ modulo eight.  Elementary
sum-to-product identities then factor this numerator into Chebyshev factors
and fixed real combinations of consecutive Chebyshev polynomials.  Such a
combination has only simple real zeros by strict interlacing.  The known
linear factor $t+1/\rt$ can therefore be removed without losing
real-rootedness.  Endpoint signs locate the unique exceptional zero below
$-1$.

\section{Factorization and real-rootedness}

Put
\[
 a=-\frac1{\rt},\qquad
 R_j(t)=U_j(t)+\rt U_{j-1}(t),\qquad
 S_j(t)=T_j(t)+\rt T_{j-1}(t).
\]

\begin{lemma}\label{lem:combination}
For every $j\geq1$, both $R_j$ and $S_j$ have $j$ distinct real zeros.
All zeros lie below $1$.  The polynomial $R_j$ has one zero below $-1$
exactly when $j\geq3$, whereas $S_j$ always has exactly one zero below $-1$.
All remaining zeros lie in $(-1,1)$.
\end{lemma}

\begin{proof}
The zeros of consecutive members of either Chebyshev family are simple and
strictly interlace.  More generally, if $P_j,P_{j-1}$ are consecutive
members of a real orthogonal-polynomial family with positive leading
coefficients, then $P_j+cP_{j-1}$ has one zero in every component of
$\R\setminus Z(P_{j-1})$: at each zero of $P_{j-1}$ its sign is the sign of
$P_j$, and those signs alternate, including against the signs at
$\pm\infty$.  This supplies $j$ distinct real roots and exhausts the degree.
Apply this with $c=\rt$ to $U_j,U_{j-1}$ and $T_j,T_{j-1}$.

At $1$, all four Chebyshev values occurring here are positive, so the largest
root is below $1$.  At the other endpoint,
\[
 R_j(-1)=(-1)^j\bigl(j+1-\rt j\bigr),\qquad
 S_j(-1)=(-1)^j(1-\rt).
\]
The sign at $-\infty$ is $(-1)^j$.  Thus the leftmost root of $R_j$
crosses below $-1$ precisely when $j+1-\rt j<0$, i.e. $j\geq3$; the
leftmost root of $S_j$ is always below $-1$.
\end{proof}

\begin{theorem}\label{thm:main}
For every $n\geq0$, all zeros of $p_n$ are real and simple.  More precisely,
write $n=2m-1$ when $n$ is odd and $n=2m$ when it is even.  Then
\begin{equation}\label{eq:numerator}
 2(t-a)q_n(t)=U_{n+1}(t)+\rt U_n(t)+c_n,
 \qquad c_n=\rt\sin\frac{3\pi n}{4},
\end{equation}
and the numerator has the following factorization:
\begin{center}
\begin{tabular}{ccl}
\toprule
condition & $c_n$ & $2(t-a)q_n(t)$ \\
\midrule
$n\equiv0\pmod4$ & $0$ & $R_{n+1}(t)$ \\
$n=2m-1$, $c_n=1$ & $1$ & $2T_m(t)R_m(t)$ \\
$n=2m-1$, $c_n=-1$ & $-1$ & $2U_{m-1}(t)S_{m+1}(t)$ \\
$n=2m$, $c_n=-\rt$ & $-\rt$ & $2T_{m+1}(t)R_m(t)$ \\
$n=2m$, $c_n=\rt$ & $\rt$ & $2U_m(t)S_{m+1}(t)$ \\
\bottomrule
\end{tabular}
\end{center}
For $n=1,2,3$, all $n$ zeros of $q_n$ lie in $(-1,1)$.  For every
$n\geq4$, exactly one zero lies in $(-\infty,-1)$ and the remaining $n-1$
lie in $(-1,1)$.  Consequently, for $n\geq4$, $p_n$ has exactly one zero
below $-1-\rt$ and all other zeros in $(-1-\rt,-1+\rt)$.  Finally,
$\lambda=-2$ is a zero of $p_n$ exactly when $n\equiv2\pmod4$.
\end{theorem}

\begin{proof}
The identity $T_n+tU_{n-1}=U_n$ and the recurrence
$U_{n+1}=2tU_n-U_{n-1}$ turn \eqref{eq:source} into
\eqref{eq:numerator}.  Since $a=\cos(3\pi/4)$,
\[
 U_{n-1}(a)=\frac{\sin(3\pi n/4)}{\sin(3\pi/4)}
 =\rt\sin\frac{3\pi n}{4};
\]
the values for $n=0,\ldots,7$ modulo eight are
$0,1,-\rt,1,0,-1,\rt,-1$.

The displayed factorization table follows from the standard product
identities
\[
\begin{aligned}
2T_mU_m&=U_{2m}+1,&2T_mU_{m-1}&=U_{2m-1},\\
2U_{m-1}T_{m+1}&=U_{2m}-1,&2U_{m-1}T_m&=U_{2m-1},\\
2T_{m+1}U_m&=U_{2m+1},&2T_{m+1}U_{m-1}&=U_{2m}-1,\\
2U_mT_{m+1}&=U_{2m+1},&2U_mT_m&=U_{2m}+1.
\end{aligned}
\]
Every factor is real-rooted by Lemma~\ref{lem:combination} and the usual
Chebyshev zero formulas.  The numerator always vanishes at $a$, so division
by $t-a$ proves real-rootedness of $q_n$.

It remains to check multiplicities.  A zero common to $T_m$ and $R_m$, or to
$U_{m-1}$ and $S_{m+1}$, would force $\cos\theta=-\rt$, which is impossible.
In the two even rows, the same trigonometric calculation shows that a common
zero can only be $a=\cos(3\pi/4)$; it occurs exactly when
$m\equiv1$ or $3\pmod4$, respectively, equivalently
$n\equiv2\pmod4$.  Thus the numerator has a double zero at $a$ in precisely
that congruence class, and division leaves one simple zero there.  When
$n\equiv0\pmod4$, Lemma~\ref{lem:combination} says the root $a$ of
$R_{n+1}$ is simple and division removes it.  All other numerator roots are
simple and belong to only one factor.  This proves simplicity and the final
criterion $\lambda=\rt a-1=-2$.

Finally, Lemma~\ref{lem:combination} shows that the $R$ or $S$ factor has
exactly one root below $-1$ in every case with $n\geq4$, while every ordinary
$T$ or $U$ factor has its roots in $(-1,1)$.  The removable root $a$ lies in
$(-1,1)$.  Directly, $q_1,q_2,q_3$ have all roots in $(-1,1)$.  The affine
map $\lambda=\rt t-1$ gives the stated intervals for $p_n$.
\end{proof}

\section{Verification and novelty bounds}

The accompanying script checks the table by exact symbolic arithmetic for
$0\leq n\leq80$ and checks real-rootedness, simplicity, and endpoint counts
for the consecutive combinations through degree $200$.  The command
\begin{quote}\small\ttfamily
conda run --no-capture-output -n sandbox python\\
code/verify\_factorizations.py
\end{quote}
returns a pass for 81 exact factorizations and for the real-root/location
checks through degree 200.  These are regression checks; the proof above is
exact.

A bounded novelty search on 21 August 2026 covered the four run indexes,
exact arXiv-id and title searches, the source question's exact zero-distribution
phrases, and later primary work on Jacobi-type pencils that cites the source.
No later paper stating this factorization or resolving this exact zero problem
was located.  Novelty therefore remains subject to expert literature review.

\paragraph{Human review recommendation.}
Check the five factorization rows, the common-zero calculation at
$a=-1/\rt$, and the endpoint sign count in Lemma~\ref{lem:combination}.

\begin{thebibliography}{9}
\bibitem{Z}
S. M. Zagorodnyuk,
\emph{Orthogonal polynomials related to some Jacobi-type pencils},
Ukrainian Math. J. \textbf{68} (2016), 1353--1365;
arXiv:1508.01794.
\bibitem{Szego}
G. Szeg\H{o}, \emph{Orthogonal Polynomials}, 4th ed.,
American Mathematical Society, 1975.
\end{thebibliography}

\end{document}
